\section{Finite witness certificates}\label{sec:certificates}

Before paying for another OBBT round, a solver would like to know how much
that round, or all remaining rounds, can still change. This section bounds
the change from above with finite sets of points that are feasible for the
relaxation. The bounds concern a specified tightening procedure and a
specified relaxation family. They are not new variable bounds, and they do
not by themselves predict solving time. Three statements must be kept apart.
\begin{enumerate}
\item A point feasible for the current relaxation bounds the improvement
  of one frozen round (Section~\ref{sec:cert-current}).
\item Points feasible for the relaxation rebuilt on an inner box $P$ bound
  the improvement of every later round, directional update, and compatible
  cutoff, provided the box being tightened contains $P$
  (Sections~\ref{sec:cert-protected}--\ref{sec:cert-scope}).
\item A residual and a verified sensitivity bound limit the remaining
  movement of an exact iteration (Section~\ref{sec:residual}).
\end{enumerate}
Statement~1 is the feasibility filtering of Gleixner et
al.~\cite{gleixner2017-three-enhancements-for-optimization-based}.
Statement~2 rests on the persistence of fixed boxes
(Lemma~\ref{lem:fixed-persist}), the order argument behind monotone
domain-reduction iterations
\cite{tarski1955fixpoint,caprara2010-global-optimization-problems-and-domain}.
The OBBT-specific content of this section is a finite certificate for a
fixed box, checked exactly on the rebuilt relaxation; a description of the
operations it covers and of those it does not; its exactness at the limit
of the iteration; and the difference between finding such a certificate
and checking it.

\subsection{Witnesses and lifted reuse}\label{sec:cert-setting}

We use the setting of Section~\ref{sec:foundations}: the relaxation family
is valid~\eqref{eq:validity} and monotone~\eqref{eq:monotonicity}, each
$\phi_B$ is lower semicontinuous, and a lifted family has compact sets
$R(B)$ and a continuous objective $v$. The cutoff sets $K_U(B)$, the
operator $T_U$, the relaxation bound $L(B)$, fixed boxes, and the Jacobi,
directional, sequential, and selective update schemes are as defined
there.

A \emph{lifted witness} on a box $P$ at cutoff $U$ is a complete lifted
vector $z\in R(P)$, including every auxiliary coordinate, that satisfies
every row of the relaxation rebuilt on $P$ and has $v(z)\le U$. By
\eqref{eq:projected}, its projection $\pi(z)$ satisfies
$\phi_P(\pi(z))\le v(z)\le U$, so $\pi(z)\in K_U(P)$. This direction needs
no attainment premise; a checked lifted witness always certifies a
projected one.

Two order properties play different roles below. All coordinate and
objective-value ceilings in this section need only the monotonicity
\eqref{eq:monotonicity} of the projected objectives. Lifted monotonicity,
$R(B')\subseteq R(B)$ for $B'\subseteq B$ with a common objective, is
needed for a different use: to treat a stored lifted point $z\in R(P)$
itself, with its stored value $v(z)$, as a feasible point of the linear
program on a box $B\supseteq P$. Such reuse lets a solver screen a later
frozen round, mix witnesses for a new cutoff (Section~\ref{sec:cutoff}),
or warm-start a later solve without recomputing auxiliary coordinates.
Under monotonicity alone, $\phi_B(\pi(z))\le\phi_P(\pi(z))\le v(z)$ still
holds, but a feasible lift on $B$ must be recomputed.

\paragraph{The reference family.}
Several examples use the lifted linear family with McCormick product rows
listed in Section~\ref{sec:foundations}. Its rows are of three kinds. Fixed rows
$Gz\le g$ do not depend on the box; they include the original linear
constraints, quadratic constraints written linearly in lifted product
variables, and valid inequalities such as $y_k\ge0$ for a lifted square.
Bound rows are $\ell\le x\le u$. For each lifted variable $y_k$ standing
for a product $x_ix_j$, with $i=j$ allowed, the four McCormick
inequalities~\cite{mccormick1976-computability-of-global-solutions-to} on
$B=[\ell,u]$ can be written as
\begin{equation}\label{eq:mccormick-gap}
\begin{aligned}
 y_k-x_ix_j&\ge-\min\{(x_i-\ell_i)(x_j-\ell_j),\,(u_i-x_i)(u_j-x_j)\},\\
 y_k-x_ix_j&\le\min\{(u_i-x_i)(x_j-\ell_j),\,(x_i-\ell_i)(u_j-x_j)\}.
\end{aligned}
\end{equation}
After the product $x_ix_j$ cancels, each of the four inequalities is
affine in $(x,y_k)$. For example, the first reads
$y_k\ge\ell_jx_i+\ell_ix_j-\ell_i\ell_j$. When $i=j$ the two upper
inequalities coincide, and the rows are the tangents of $x_i^2$ at both
endpoints together with the secant. This is a finite linear relaxation of
the square, not its convex epigraph. The objective is a fixed affine
function $v(z)=c^Tz+c_0$.

\begin{lemma}[The reference family]\label{lem:reference-family}
The reference family is lifted monotone with a common objective, and hence
monotone. Each $R(B)$ is a polytope. If the fixed rows are valid for the
original problem and $v$ evaluated at the graph lift $y_k=x_ix_j$ equals
$f(x)$, the family is valid.
\end{lemma}
\begin{proof}
Let $B'=[\ell',u']\subseteq B=[\ell,u]$ and $x\in B'$. Then
$0\le x_i-\ell'_i\le x_i-\ell_i$ and $0\le u'_i-x_i\le u_i-x_i$ for every
$i$. Each product in~\eqref{eq:mccormick-gap} multiplies two such
nonnegative distances, so its value on $B'$ is at most its value on $B$.
Consequently each lower bound on $y_k-x_ix_j$ on $B'$ is at least the
corresponding lower bound on $B$, and each upper bound is at most the
corresponding upper bound. Fixed rows and the objective do not depend on
the box, and bound rows on $B'$ imply those on $B$. Thus
$R(B')\subseteq R(B)$. The set $R(B)$ is a polyhedron; it is bounded
because $x\in B$ and each $y_k$ lies between affine functions of $x$.
At the graph lift the middle terms in~\eqref{eq:mccormick-gap} vanish,
and all bounds have the correct sign because every distance is
nonnegative on $B$. Together with valid fixed rows, this puts the graph
lift of each $x\in X\cap B$ in $R(B)$, so $\phi_B(x)\le f(x)$.
\end{proof}

\begin{example}[Monotonicity without lifted monotonicity]\label{ex:lambda}
Relax $s=x^2$ on a box $[\ell,u]$ by the convex-combination formulation
common in piecewise-linear relaxations: auxiliary weights
$\lambda_0,\lambda_1\ge0$ with $\lambda_0+\lambda_1=1$ and
$x=\lambda_0\ell+\lambda_1u$, the secant row
$s\le\lambda_0\ell^2+\lambda_1u^2$, the two endpoint tangents, and the
objective $v=s$. In every case
$\lambda_0\ell^2+\lambda_1u^2=(\ell+u)x-\ell u$, and for $\ell<u$ the
weights are determined by $x$. The projection onto
$(x,s)$ is therefore the square relaxation of the reference family, and
since the objective depends only on $s$, the projected objectives are
those of that family, which are monotone by
Lemma~\ref{lem:reference-family}. Lifted monotonicity fails: on
$P=[0,1/2]$ the point $x=1/2$ has weights $(0,1)$, while on $B=[0,1]$ it
requires $(1/2,1/2)$. A stored lifted witness for $P$ is not a feasible
point of the relaxation on $B$, although its projection remains in every
projected set that the certificates below use.
\end{example}

\subsection{The current frozen round}\label{sec:cert-current}

\begin{proposition}[Current-round ceilings]\label{prop:current-round}
Let $B=[\ell,u]$, let $W$ be a nonempty finite subset of $K_U(B)$, and
write $T_U(B)=[\ell',u']$. For every coordinate $i$,
\begin{equation}\label{eq:current-ceilings}
 0\le\ell'_i-\ell_i\le\min_{x\in W}x_i-\ell_i,\qquad
 0\le u_i-u'_i\le u_i-\max_{x\in W}x_i.
\end{equation}
\end{proposition}
\begin{proof}
The coordinate minimum over $K_U(B)$ is at most the minimum over its
subset $W$, and the maximum is at least the maximum over $W$. Since
$K_U(B)\subseteq B$, the new endpoints lie in $[\ell_i,u_i]$.
\end{proof}

A witness attaining an endpoint proves that this direction cannot move in
the present round, and one point can screen several directions. The
ceilings bound possible benefit; a witness coordinate is not a valid new
bound. Points produced by earlier support solves of the same round, or
an optimal solution of the relaxation itself, are typical witnesses.
Checking $k$ lifted points against $N$ nonzero row coefficients costs
$O(kN)$ arithmetic operations. Finding good points can cost as much as the
round itself. A ceiling on coordinate movement also says nothing about
side products of a support solve, such as dual information used for
Lagrangian variable
bounds~\cite{gleixner2017-three-enhancements-for-optimization-based}.

The ceiling~\eqref{eq:current-ceilings} concerns the frozen relaxation
on $B$. After rebuilding on a smaller box, the same point may be
infeasible, and later rounds may move much further.

\begin{example}[A frozen witness does not bound later rounds]\label{ex:halving}
Minimize $f(x)=x^2$ over $B_0=[-r_0,r_0]$, $r_0>0$, with the square
relaxation of the reference family, objective $v=s$, and the cutoff $U=0$
given by the incumbent $x=0$; this is the case $n=1$ of
Example~\ref{ex:finite-square}. On $[-r,r]$ the rows are $s\ge2r|x|-r^2$
(the two tangents) and $s\le r^2$ (the secant). With $s\le0$ they give
$|x|\le r/2$, and every such $x$ has the lift $s=0$. Hence
$T_0([-r,r])=[-r/2,r/2]$, and the exact Jacobi iterates are
$B_k=[-2^{-k}r_0,2^{-k}r_0]$. In the first round the witness $(r_0/2,0)$
bounds the decrease of the upper endpoint by $r_0/2$, which is exact for
that round. After rebuilding on $[-r_0/2,r_0/2]$, the new upper tangent
requires $s\ge r_0x-r_0^2/4=r_0^2/4>0$ at $x=r_0/2$, so the witness is
excluded. The total decrease over all rounds is $r_0$.
\end{example}

\subsection{Protected boxes}\label{sec:cert-protected}

A witness remains useful in all later rounds if it is feasible for the
relaxation built on a box that every later box contains. The following
definitions make this precise and describe the operations that preserve
such a box.

\begin{definition}[Relaxation-sound step]
A box $C'$ is obtained from a box $C$ by a relaxation-sound step at
cutoff $U'$ if $T_{U'}(C)\subseteq C'\subseteq C$.
\end{definition}

A Jacobi round is a relaxation-sound step with $C'=T_{U'}(C)$. So is a
directional update, and so is the simultaneous update of any selected set
of endpoints on one frozen relaxation. Sequential and selective rounds
are sequences of such steps. Conservative versions, which move each
selected endpoint at most to its exact support, are also relaxation-sound.
This includes OBBT whose endpoint bounds are proved by safe dual bounds
evaluated with directed rounding (Proposition~\ref{prop:dual-residual}).
Rounding an approximate optimum outward alone does not prove such a bound.
The identity step is relaxation-sound.

\begin{definition}[Witness pool and protected box]
Let $P=[a,b]$ be a nonempty box. A \emph{witness pool} for $P$ is a
finite set $W\subseteq K_U(P)$ with $\hull W=P$; that is, for every
coordinate $i$ some $x\in W$ has $x_i=a_i$ and some $x'\in W$ has
$x'_i=b_i$. The \emph{threshold} of $W$ is
$U_W=\max_{x\in W}\phi_P(x)\le U$. A box with a witness pool is
\emph{protected} at every cutoff $U'\ge U_W$.
\end{definition}

The definition evaluates $\phi_P$, the relaxation \emph{rebuilt on $P$},
not the relaxation on a larger current box. In lifted form, a pool
consists of lifted witnesses on $P$; their stored values $v(z)$ bound the
threshold from above. At most $2n$ points are needed, one for each
endpoint, and fewer suffice when one point attains several endpoints. A
protected box is a fixed box; the word \emph{protected} refers to the
finite certificate and its threshold.

\begin{theorem}[Protection]\label{thm:protected}
Let $W$ be a witness pool for $P=[a,b]$ with threshold $U_W$. Let
$C_0\supseteq P$ be a box in $B_0$, and for $k\ge0$ let $C_{k+1}$ be
obtained from $C_k$ by a relaxation-sound step at a cutoff $U_k\ge U_W$.
Then:
\begin{enumerate}[label=(\alph*)]
\item $T_{U'}(P)=P$ for every $U'\ge U_W$;
\item $P\subseteq C_k$ and $W\subseteq K_{U_k}(C_k)$ for every $k$; in
  particular no step proves $C_k$ infeasible at its cutoff;
\item writing $C_k=[\ell^k,u^k]$, for every $i$ and $k$,
\begin{equation}\label{eq:protected-ceilings}
 0\le\ell^k_i-\ell^0_i\le a_i-\ell^0_i,\qquad
 0\le u^0_i-u^k_i\le u^0_i-b_i.
\end{equation}
\end{enumerate}
\end{theorem}
\begin{proof}
(a) For $U'\ge U_W$, $W\subseteq K_{U_W}(P)\subseteq K_{U'}(P)$, so
$P=\hull W\subseteq T_{U'}(P)\subseteq P$.
(b) By (a), $P$ is a fixed box of $T_{U_W}$. If $P\subseteq C_k$, then
Lemma~\ref{lem:fixed-persist} gives $P\subseteq T_{U_k}(C_k)\subseteq
C_{k+1}$, and induction gives $P\subseteq C_k$ for every $k$.
Lemma~\ref{lem:order} gives $W\subseteq K_{U_W}(P)\subseteq
K_{U_k}(C_k)$.
(c) The inequalities restate $P\subseteq C_k\subseteq C_0$.
\end{proof}

The theorem needs no continuity, convergence, or fairness in the choice of
directions; the steps may mix Jacobi rounds and directional updates and
may use different cutoffs above $U_W$. The witnesses need not be
feasible for the original problem, and none of them needs to be an
optimal solution of a support problem. The essential requirement is
feasibility in the relaxation rebuilt on $P$. In
Example~\ref{ex:halving} the first-round witnesses are feasible on $B_0$
but not on their own hull, and indeed that hull is not retained.

A protected box limits what the specified OBBT procedure can remove. It
is not a domain reduction and never replaces the search domain. For
instance, minimize $f(x)=x$ over $[0,1]$ with the exact relaxation
$\phi_B(x)=x$ and the cutoff $U=1$ of the incumbent $\hat x=1$. The
singleton $\{1\}$ is protected by the witness $\hat x$, and every OBBT
iterate contains it, but the minimizer $0$ lies outside it. Optimal
solutions may lie anywhere in the current box. A protected box need not
contain any original feasible point either: its pool may consist of
relaxation-only points.

The finite certificate loses nothing when a prescribed box is tested: a
box is fixed exactly when a short witness pool exists.

\begin{proposition}[Fixed boxes have short certificates]\label{prop:fixed-box}
A box with a witness pool is a fixed box. Conversely, every fixed box $P$
of $T_U$ has a witness pool $W\subseteq K_U(P)$ with $|W|\le2n$, and for
a lifted family the pool can be chosen to consist of at most $2n$ lifted
witnesses. If, moreover, $R(P)$ is a polytope given by rational data, $v$
has rational coefficients, and the endpoints of $P$ and $U$ are rational,
then these lifted witnesses can be chosen rational, so an exact rational
check accepts them.
\end{proposition}
\begin{proof}
The first claim is Theorem~\ref{thm:protected}(a). Let $P$ be fixed. The
set $K_U(P)$ is a sublevel set of the lower semicontinuous function
$\phi_P$ on the compact box $P$, hence compact, and $T_U(P)=P$ means that
each coordinate attains its minimum $a_i$ and maximum $b_i$ on it.
Choosing one minimizer and one maximizer per coordinate gives at most $2n$
points whose hull is $P$. For a lifted family, the infimum
in~\eqref{eq:projected} is attained, so each chosen point has a lift in
$R(P)$ with value at most $U$. In the rational case, the set of lifts
attaining a given endpoint, such as $\{z\in R(P):v(z)\le U,\ x_i=a_i\}$,
is a nonempty polytope defined by rational data, and any of its vertices
is the solution of a nonsingular rational linear system.
\end{proof}

The proposition makes the certificate complete as a \emph{test} for a
prescribed box. It is not a search procedure: it says nothing about how
to find a fixed box, and Section~\ref{sec:cert-cost} shows that natural
proposals can miss one.

Certificates from different sources can be combined.

\begin{proposition}[Combining certificates]\label{prop:join}
If $W_1$ and $W_2$ are witness pools for $P_1$ and $P_2$ with thresholds
$U_1$ and $U_2$, then $W_1\cup W_2$ is a witness pool for
$\hull(P_1\cup P_2)$ with threshold at most $\max\{U_1,U_2\}$.
\end{proposition}
\begin{proof}
Let $P=\hull(P_1\cup P_2)$ and $U'=\max\{U_1,U_2\}$. For $x\in W_j$,
$\phi_P(x)\le\phi_{P_j}(x)\le U_j\le U'$ by monotonicity. Every endpoint
of $P$ is an endpoint of $P_1$ or $P_2$ and is attained by the
corresponding pool.
\end{proof}

For example, the singleton of an incumbent that lies in the current box,
which is protected (Corollary~\ref{cor:original-points}), can be merged
with a stalled box that does not contain the incumbent. The ceilings of the merged box are
no larger than those of either part, and some are smaller.

\subsection{Exact ceilings at the limit}\label{sec:cert-limits}

Theorem~\ref{thm:protected} gives upper bounds on remaining movement. The
best of these bounds is exact whenever the limit of the iteration is a
fixed box.

\begin{corollary}[Exact ceilings at the limit]\label{cor:limit-ceilings}
Let $U\ge f^*$, let $B_k$ be the Jacobi iterates from $B_0$, and assume
the closedness condition~\eqref{eq:closed-family}. Then the limit
$B_\infty$ has a witness pool of at most $2n$ points. For every $k$, the
ceilings~\eqref{eq:protected-ceilings} of $P=B_\infty$ from $C_0=B_k$
equal the exact remaining movements of the Jacobi iteration. The same
holds for complete sequential rounds $S_{k+1}=S_U(S_k)$ with $S_0=B_0$:
$B_\infty\subseteq S_k$, the iterates $S_k$ converge to $B_\infty$, and
the ceilings of $B_\infty$ from $S_k$ equal their remaining movements.
\end{corollary}
\begin{proof}
By Proposition~\ref{prop:fixed-limit}, $B_\infty$ is a fixed box, and
Proposition~\ref{prop:fixed-box} gives the pool. The ceilings of
$B_\infty$ from $B_k$ are the distances between corresponding endpoints
of $B_k$ and $B_\infty$, which are the remaining movements because $B_k$
converges to $B_\infty$. A complete sequential round is a sequence of
directional updates, so Theorem~\ref{thm:protected} gives
$B_\infty\subseteq S_k$, and Lemma~\ref{lem:sequential} shows that the
$S_k$ have the same intersection $B_\infty$.
\end{proof}

The reference family satisfies~\eqref{eq:closed-family} by the argument
following Proposition~\ref{prop:fixed-limit}, with the compact set
$R(B_0)$ containing every $R(B)$ by Lemma~\ref{lem:reference-family}.
Under this condition, no information is lost by measuring remaining
benefit through protected boxes: the exact remaining movement is a
protected-box ceiling. The difficulty lies entirely in finding the right
box. Without~\eqref{eq:closed-family} the corollary can fail. In
Example~\ref{ex:not-fixed} the greatest fixed box is $\{0\}$, while the
iterates converge to $[0,1/2]$; from $B_k=[0,s_k]$ the ceiling of
$\{0\}$ is $s_k$, whereas the remaining movement is $s_k-1/2$.

\begin{example}[Exact ceilings without finite termination]\label{ex:sharp-halving}
Return to Example~\ref{ex:halving}. The limit is $B_\infty=\{0\}$. It is
protected by the lifted witness $(x,s)=(0,0)$; on the degenerate box the
square rows force $s=0$. The ceilings~\eqref{eq:protected-ceilings} from
$B_0$ are $r_0$ in both directions, and they equal the total movement,
although no finite number of rounds reaches $\{0\}$. The relaxation bound
behaves in the same way: $L(B_k)=-4^{-k}r_0^2$, attained at
$(0,-4^{-k}r_0^2)$, while Corollary~\ref{cor:objective-ceiling} below
gives the ceiling $0$, which is the limit. The witnesses produced by any
single round, $(\pm r/2,0)$ on $[-r,r]$, never pass the rebuilt check.
\end{example}

\subsection{Objective ceilings}\label{sec:cert-objective}

The relaxation bound $L(C)=\inf_{x\in C}\phi_C(x)$ of
Section~\ref{sec:foundations} is also controlled by a protected box.

\begin{corollary}[Objective ceiling]\label{cor:objective-ceiling}
Let $W$, $P$ and $(C_k)$ be as in Theorem~\ref{thm:protected}. For every
$k$ and every $x\in P$,
\begin{equation}\label{eq:objective-ceiling}
 L(C_0)\le L(C_k)\le L(P)\le\phi_P(x).
\end{equation}
In particular, a lifted point $z\in R(P)$ of any objective value gives
$L(C_k)\le v(z)$. If $\lambda\le L(C_0)$, the relaxation bound can
increase by at most $v(z)-\lambda$ along the sequence.
\end{corollary}
\begin{proof}
Since $C_k\subseteq C_0$, monotonicity of $L$ gives $L(C_k)\ge L(C_0)$.
Since $P\subseteq C_k$ and $\phi_{C_k}\le\phi_P$ on $P$,
$L(C_k)\le\inf_{x\in P}\phi_P(x)=L(P)$.
\end{proof}

The objective point need not be feasible for the original problem, and
it need not be optimal for the relaxation on $P$; any lifted point
feasible on $P$ gives a valid ceiling, and the minimum of $v$ over $R(P)$
gives the best one. Monotonicity suffices for this value statement.
Under lifted monotonicity the simpler argument applies: $z$ itself stays
feasible in every $R(C_k)$ with the same value. The lower bound $\lambda$
must bound the value of \emph{this} relaxation on $C_0$. A lower bound on
the original optimum cannot be substituted, because the relaxation value
can lie far below it; Example~\ref{ex:three-variable} has $L=-3$ and
original optimum $0$, and substituting $0$ would give a negative ceiling.

The practical reading concerns pruning. If $L(P)<U$, then
$U-L(C_k)\ge U-L(P)>0$ for every $k$: no amount of OBBT with this family
closes the gap at this node. Branching, cuts, or a stronger relaxation
are needed. The corollary bounds only the relaxation value along a
sequence of relaxation-sound steps. A child node created by branching need
not contain $P$, and a stronger relaxation has a different value.

\begin{example}[A strongly convex objective with a permanent gap]\label{ex:three-variable}
Minimize
\[
 f(x)=x_1^2+x_2^2+x_3^2+x_1x_2+x_1x_3+x_2x_3
     =\tfrac12\bigl(\norm{x}_2^2+(x_1+x_2+x_3)^2\bigr)
\]
over $[-1,1]^3$, the case $a=1$, $n=3$ of Example~\ref{ex:many-term}. The
unique minimizer is $0$, with $f^*=0$. Use the reference family with
square variables $s_i$, product variables $y_{ij}$ for $i<j$, fixed rows
$s_i\ge0$, objective $v=\sum_is_i+\sum_{i<j}y_{ij}$, and the cutoff $U=0$
of the optimal incumbent. On the cube $P=[-1,1]^3$ the square rows are
$s_i\ge2|x_i|-1$ and $s_i\le1$, and the product rows are
$|x_i+x_j|-1\le y_{ij}\le1-|x_i-x_j|$.

For the endpoint $x_1=1$, take $x=(1,0,0)$, $s=(1,0,0)$,
$y_{12}=y_{13}=0$ and $y_{23}=-1$. Every row holds: $s_1=1$ meets its
tangent and secant, the other squares satisfy $0\ge-1$, the incident
products satisfy $0\le y_{1j}\le0$, and the opposite product satisfies
$-1\le y_{23}\le1$. Its value is $1-1=0\le U$. The symmetric points for
the other five endpoints have the same value. Hence the cube is protected
with threshold $0=f^*$. Every relaxation-sound step at every cutoff
produced by a feasible incumbent leaves the whole cube unchanged.

The point $x=0$, $s=0$, $y_{ij}=-1$ has value $-3$, so $L(C_k)\le-3$ for
every such sequence. The value is exactly $-3$: $s_i\ge0$ and
$y_{ij}\ge|x_i+x_j|-1\ge-1$. The gap $U-L=3$ therefore persists. All
these points satisfy $s_i\ge x_i^2$; replacing the finite square rows by
the convex epigraphs would not remove the stall, in agreement with the
tangent-model analysis of Example~\ref{ex:many-term}. The cause is the
accumulated McCormick gap of the three products. This is the equality
case of condition~(a) of Proposition~\ref{prop:quadratic-rows}:
$H_{kk}/2=S_k=1$ for every $k$.
\end{example}

\subsection{Original feasible points}\label{sec:cert-original}

Validity gives protected boxes without any relaxation solve.
Lemma~\ref{lem:sublevel-hull} shows that every iterate contains the hull
$H_U$ of the whole original sublevel set. A finite part of that set gives
a certificate that a checker can verify.

\begin{corollary}[Hulls of original feasible points]\label{cor:original-points}
Let $S$ be a finite nonempty subset of $X\cap B_0$ with $f(x)\le U$ for
all $x\in S$, and let $P=\hull S$. Then $S$ is a witness pool for $P$
with threshold at most $\max_{x\in S}f(x)$. Moreover, $P$ is retained by
every box operation that keeps all points of
$\{x\in X\cap C:f(x)\le U'\}$ in the current box $C$, for cutoffs
$U'\ge\max_{x\in S}f(x)$.
\end{corollary}
\begin{proof}
Each $x\in S$ lies in $X\cap P$, so validity gives $\phi_P(x)\le f(x)\le
U$, and $\hull S=P$. An operation that keeps every original feasible
point with objective at most $U'$ keeps $S$, hence its hull.
\end{proof}

The second statement covers more operations than
Theorem~\ref{thm:protected}: every reduction that is valid for the
original problem, including cuts, propagation of original constraints,
stronger relaxations, and integer rounding when $S$ respects integrality.
In that case the integer-coordinate endpoints of $P$ are automatically
integral. The singleton of the incumbent is the simplest instance: OBBT
never removes the incumbent from a box that contains it. At a node whose
box excludes the incumbent, the incumbent still supplies the cutoff but
no protected singleton. With several known solutions, for example
images of the incumbent under a symmetry of the model, the hull can be
large. A lifted checker needs exact lifts of the points in $S$, which can
be hard to obtain for nonlinear equalities; a numerically accepted
incumbent is not exactly feasible in general.

\subsection{Rounding, cuts, and branching}\label{sec:cert-scope}

A solver interleaves OBBT with integer rounding, cut generation,
propagation, and branching. A certificate for a fixed relaxation family
covers these operations only under the conditions below.

\begin{proposition}[Integer rounding]\label{prop:integer-rounding}
Let $I$ be the set of integer-constrained coordinates, and let rounding
replace $\ell_i$ by $\lceil\ell_i\rceil$ and $u_i$ by $\lfloor u_i\rfloor$
for $i\in I$. If $P=[a,b]$ is protected and $a_i,b_i\in\mathbb Z$ for all
$i\in I$, then the conclusions of Theorem~\ref{thm:protected} remain true
when roundings are interleaved with the relaxation-sound steps.
\end{proposition}
\begin{proof}
Rounding a box $C\supseteq P$ keeps $P$: from $\ell_i\le a_i\in\mathbb Z$
follows $\lceil\ell_i\rceil\le a_i$, and from $u_i\ge b_i\in\mathbb Z$
follows $\lfloor u_i\rfloor\ge b_i$. The induction in
Theorem~\ref{thm:protected} therefore continues through rounding steps.
\end{proof}

Fractional faces can be removed. Minimize $x^2$ over $[-1,1]$ with the
square relaxation and the cutoff $U=1/4$. The box $[-1/2,1/2]$ is
protected by the lifted witnesses $(\pm1/2,1/4)$: on it the tangents read
$s\ge|x|-1/4$ and the secant reads $s\le1/4$. It is the limit of the exact
iteration from $[-1,1]$, whose first round gives $[-5/8,5/8]$. If $x$ is
integer, the first rounding gives $\{0\}$, removing the protected faces;
here this is correct, because $0$ is the only integer point with
$x^2\le1/4$. If the support problems themselves enforce integrality, the
operator changes: $K_U$ is replaced by its subset of points satisfying
integrality, which is again monotone, and a pool must consist of such
points. Its endpoints are then integral on $I$.

New rows change the family. If a lifted monotone family receives the same
row on every box, lifted monotonicity is preserved; if every lifted
witness of a pool satisfies the row, the pool remains a certificate for
the modified family. Otherwise a new check is required. A cut can remove
an objective witness while keeping the face witnesses. On $[-1/2,1/2]$ at
cutoff $1/4$ above, the lifted point $(0,-1/4)$ is feasible and gives
$L\le-1/4$, although $x^2\ge0$. The valid row $s\ge0$ excludes this point,
and the strengthened relaxation has bound $0$, while the face witnesses
$(\pm1/2,1/4)$ survive. Reductions valid for the original problem but
stronger than the relaxation can remove relaxation-only faces altogether.
A solver that recognizes convexity of $f$ in
Example~\ref{ex:three-variable} reduces the cube to $\{0\}$ at cutoff $0$.
Pools of original feasible points (Corollary~\ref{cor:original-points})
are guaranteed to survive every valid reduction that keeps their points.

A certificate applies at a node only if the node box contains $P$. The
intersection of $P$ with a child box need not be protected.
Example~\ref{ex:shape-change} shows the same objective tightening on a
thinner box for the tangent model; the next example shows it for a child
created by branching in the lifted family.

\begin{example}[A protected box does not pass to a child]\label{ex:child}
In Example~\ref{ex:three-variable}, branch on $x_1\ge0$, giving the child
$C=[0,1]\times[-1,1]^2=P\cap C$. On $C$ the product rows for $y_{12}$ are
$y_{12}\ge\max\{-x_1,\,x_1+x_2-1\}$ and the two upper rows, and similarly
for $y_{13}$. At the endpoint $x_2=1$, every relaxed point satisfies
$s_2\ge1$, $y_{12}\ge x_1$, $y_{13}\ge-x_1$, $y_{23}\ge x_2+x_3-1=x_3$,
$s_1\ge0$ and $s_3\ge\max\{0,-2x_3-1\}$. Hence
\[
 v\ge1+x_3+\max\{0,-2x_3-1\}\ge\tfrac12,
\]
where the last bound follows by considering $x_3\ge-1/2$ and
$x_3\le-1/2$. The value $1/2$ is attained at $x=(0,1,-1/2)$ with
$y_{23}=-1/2$ and all other lifted variables at their lower bounds. No
point of $K_0(C)$ attains $x_2=1$, so OBBT on the child tightens $u_2$.
The parent witness for this endpoint has $y_{13}=-1$, which violates the
child row $y_{13}\ge-x_1=0$.
\end{example}

\subsection{Finding a certificate versus checking it}\label{sec:cert-cost}

Checking a proposed certificate is inexpensive. Given a candidate box $P$
and a lifted pool of $k\le2n+1$ points, one builds the relaxation on $P$
and evaluates every row at every point: with $N$ nonzero row
coefficients, this is $O(kN)$ arithmetic operations, plus $O(kn)$
comparisons for endpoint attainment and $k$ objective evaluations. No
optimization problem is solved. The check must be exact, for example in
rational arithmetic. A feasibility test with a tolerance can accept a
point that violates a rebuilt row slightly, and then the conclusion that
no later round moves an endpoint is unfounded.

Finding a certificate is a different matter. By
Proposition~\ref{prop:fixed-box}, a candidate must be a fixed box of
$T_U$, and useful ceilings require a fixed box close to the limit of the
iteration. Computing such a box can require the support solves that OBBT
itself performs, and particular proposals can fail even at a fixed box.
Two concrete proposals are available. The first uses the points that an
exact round produces anyway.

\begin{proposition}[Checking a completed round]\label{prop:round-check}
Let $B'=T_U(B)$ be nonempty and put
$R_U(C)=\{z\in R(C):v(z)\le U\}$ for a box $C$. Let
$\widehat W\subseteq R_U(B)$ be a finite set of complete lifted points
whose projections $W=\pi_x(\widehat W)$ attain all $2n$ endpoints of $B'$,
where $\pi_x(x,y)=x$, such as lifted optimal solutions of the support problems.
\begin{enumerate}[label=(\alph*)]
\item If $\widehat W\subseteq R_U(B')$, then $W$ is a witness pool for $B'$, and
  Theorem~\ref{thm:protected} applies with $P=B'$.
\item If $T_U(B)=B$, the hypothesis of~(a) holds.
\item The hypothesis of~(a) can fail although $T_U(B')=B'$.
\end{enumerate}
Consequently, if exact Jacobi iterates reach a fixed box after finitely
many rounds, the test in~(a) succeeds no later than in the round that
starts at the fixed box. If they never reach one, the test never
succeeds.
\end{proposition}
\begin{proof}
(a) The lifted check gives $W\subseteq K_U(B')$, and by construction
$\hull W=B'$. (b) Then $B'=B$, so
$\widehat W\subseteq R_U(B)=R_U(B')$. (c) See Example~\ref{ex:round-check}. For the
consequence, if $B_j$ is the first fixed box, part~(b) applies to the
round that starts at $B_j$. If the test succeeded in some round, the
resulting box would be fixed by part~(a).
\end{proof}

\begin{example}[A fixed box missed by its own round]\label{ex:round-check}
Minimize $f(x)=x^2$ over $B=[0,1]$ subject to $x\le1/2$, kept as a fixed
row, with the square relaxation and the cutoff $U=1/4$ of the incumbent
$x=1/2$. On $B$ the rows are $s\ge0$, $s\ge2x-1$, $s\le x$, $x\le1/2$ and
$s\le1/4$. The lower support is $0$, attained at $(0,0)$. The upper
support is $1/2$, attained on the segment $\{1/2\}\times[0,1/4]$, with
vertices $(1/2,0)$ and $(1/2,1/4)$. Thus $B'=[0,1/2]$. On $B'$ the tangent
at $1/2$ and the secant force $s=1/4$ at $x=1/2$. The vertex $(1/2,0)$
fails the rebuilt check, while $(1/2,1/4)$ passes, so the outcome depends
on which optimal vertex the solver returns. The box $B'$ is fixed either
way; it is the hull of the original feasible points $0$ and $1/2$.
\end{example}

In exact arithmetic, a Jacobi round that returns its input box already
proves a fixed box. The rebuilt check adds three things. It can detect
the fixed box one round earlier and save that round's $2n$ support
solves. It produces a portable certificate, which also covers directional
updates, other cutoffs above its threshold, and every containing box.
Finally, it needs only primal points. A valid \emph{reduction} of a bound
needs a dual proof that no relaxed point lies beyond the new bound.
Proving the \emph{absence} of further reduction needs only primal
witnesses, which is exactly what a pool supplies.
Section~\ref{sec:algorithms} describes an exact driver that verifies each
round by rational primal--dual proofs and then applies the test of
Proposition~\ref{prop:round-check}; Theorem~\ref{thm:closure} states its
guarantees. When the iterates converge only asymptotically, as in
Example~\ref{ex:halving}, this proposal never succeeds, although
Corollary~\ref{cor:limit-ceilings} guarantees that the limit is
protected.

The second proposal uses original feasible points in the current box:
the incumbent, other solutions found by heuristics, and their images under
known symmetries (Corollary~\ref{cor:original-points}). It needs no relaxation solve, only
exact lifts, and it survives all valid reductions. Its ceilings are exact
when the hull equals the limit. In Example~\ref{ex:sharp-halving} the
incumbent singleton $\{0\}$ is the limit, and in
Example~\ref{ex:round-check} the hull of $0$ and $1/2$ is the fixed box.
In general the hull of known solutions can be far inside the limit.
Proposition~\ref{prop:join} merges certificates from both proposals.

A protected box therefore does not make the round that finds it cheaper.
Its value lies in later work that it makes unnecessary: rounds after a
certified stall, repeated tightening after cutoff changes above its
threshold, repeated tightening at nodes that contain it, and certified
bounds on the remaining relaxation gap. Whether these savings exceed the
cost of discovery depends on how often such events occur during search.
